% ===================================================================
%  TABELUL Nr. 6 — Casa pe dinăuntru. — Mobila. — Mâncarea. — Luminatul.
%  Traduction 2026 d'apres texto/io, controlee sur texto/fr, texto/en
%  et texto/es. Consigne : voir texto/ro/10-tabelul-01.tex.
%
%  Deux notes, toutes deux marquees « (*) », et deux appels : celui de
%  « babo » au bloc c1-12-1, celui de « lambrequino » au bloc c2-01-1.
%
%  LE BLOC c1-06-1 GARDE SON ECART DECLARE : l'ido y porte « (6) » la
%  ou le francais porte « (16) », et c'est le francais qui a raison.
%
%  LE PLUS GROS LEXIQUE DE LA COLONNE — cinq pieces, deux cent
%  cinquante objets — et le plus riche en mots turcs, qui sont ceux de
%  la maison : « dulap » l'armoire, « perdea » le rideau, « cearșaf »
%  le drap, « farfurie » l'assiette, « tavan » le plafond, « chibrit »
%  l'allumette, « ciorap » le bas. Sept objets domestiques, sept mots
%  venus par les Ottomans.
% ===================================================================

%%K t06-apar-1 apar
\VUsaut{6.0mm}
\VUtitre{15.0pt}{60}{TABELUL Nr. 6}
\VUsaut{1.4mm}
\VUtitre{11.4pt}{0}{Casa pe dinăuntru, Mobila, Mâncarea,}
\VUsaut{0.4mm}
\VUtitre{11.4pt}{0}{Luminatul.}
\VUsaut{2.2mm}

%%K t06-apar-2 apar
Casa este gata. Unchiul lui Ioan s-a mutat în noua sa locuință, a cărei împărțire
o cunoaștem deja. Să vedem acum cum a mobilat-o. Mai întâi vom cerceta:

%%K t06-tit-1 sub
\VUpk{10.2pt}{40}{Camera de culcare.}
\VUsaut{3.0mm}

%%K t06-c1-01-1 p
1. --- Am venit într-o clipă nepotrivită, fiindcă slujnica face curat în
\VUgras{cameră}.

%%K t06-c1-02-1 p
2. --- Pe \VUgras{lavoar} \textsuperscript{(1)} stau ici și colo lucruri felurite,
cu care se spală, se piaptănă și îndeobște se dichisesc. Sunt \VUgras{ligheanul}
\textsuperscript{(2)} și \VUgras{ulciorul} \textsuperscript{(3)}, \VUgras{peria
de păr} \textsuperscript{(4)}, \VUgras{pieptenele} \textsuperscript{(5)},
\VUgras{săpunul} \textsuperscript{(6)} și \VUgras{buretele}
\textsuperscript{(7)}, în sfârșit \VUgras{sticluța} \textsuperscript{(8)} cu apă
de gură.

%%K t06-c1-03-1 p
3. --- Pe podea, înaintea lavoarului, stă o \VUgras{găleată}
\textsuperscript{(9)} pentru apa murdară.

%%K t06-c1-04-1 p
4. --- În \VUgras{antreu} \textsuperscript{(10)} se văd câteva \VUgras{cuiere}
\textsuperscript{(13)} și două \VUgras{umbrele} \textsuperscript{(11)} într-un
\VUgras{suport} \textsuperscript{(12)}.

%%K t06-c1-05-1 p
5. --- De cealaltă parte a ușii stă \VUgras{dulapul cu oglindă}
\textsuperscript{(14)}, în care se păstrează \VUgras{rufăria}
\textsuperscript{(15)}.

%%K t06-c1-06-1 p
6. --- Să folosim faptul că \VUgras{camerista} \textsuperscript{(16)} face
\VUgras{patul} \textsuperscript{(17)}, și să îi cercetăm părțile. \VUgras{Patul}
\textsuperscript{(20)} poartă o \VUgras{somieră} bună \textsuperscript{(21)}, pe
care Justina va așeza îndată \VUgras{salteaua} de lână \textsuperscript{(22)}.
Apoi va lua de pe scaune două \VUgras{cearșafuri} \textsuperscript{(23)} și
\VUgras{pătura} \textsuperscript{(24)}. Pe urmă va pune la căpătâi
\VUgras{sulul} \textsuperscript{(25)} și \VUgras{perna} \textsuperscript{(26)},
care stau împreună cu \VUgras{plapuma de puf} \textsuperscript{(27)} pe
\VUgras{covorul de lângă pat} \textsuperscript{(32)}. Cu pămătuful va scutura
\VUgras{perdelele} \textsuperscript{(19)} și \VUgras{baldachinul}
\textsuperscript{(18)}, și o parte din treabă va fi făcută.

%%K t06-c1-07-1 p
7. --- Apoi va trebui să facă puțină rânduială pe \VUgras{noptieră}
\textsuperscript{(28)}, să întoarcă \VUgras{deșteptătorul}
\textsuperscript{(29)}, să pregătească \VUgras{lumânarea de noapte}
\textsuperscript{(30)} și să ducă \VUgras{lumânarea} \textsuperscript{(31)}, ca să
curețe sfeșnicul.

%%K t06-tit-2 sub
\VUsaut{5.0mm}
\VUpk{10.2pt}{40}{Coșul de lucru și garnitura căminului.}
\VUsaut{3.0mm}

%%K t06-c1-08-1 p
8. --- \VUgras{Coșul de lucru} \textsuperscript{(34)} de lângă \VUgras{taburet}
\textsuperscript{(33)} are și el mare nevoie de rânduială. Va trebui împăturită
\VUgras{broderia} \textsuperscript{(35)}, puse în \VUgras{măsuța de lucru}
\textsuperscript{(38)} \VUgras{degetarul}, \VUgras{ațele}
\textsuperscript{(36)}, \VUgras{acele} \textsuperscript{(37)} și
\VUgras{foarfecele} \textsuperscript{(39)}, și încuiat capacul măsuței cu
\VUgras{cheia} \textsuperscript{(40)}.

%%K t06-c1-09-1 p
9. --- Pe \VUgras{măsuța rotundă} \textsuperscript{(41)} din mijlocul camerei
Justina va schimba apa din \VUgras{carafă} \textsuperscript{(42)}. Va pune
\VUgras{zahărul} în \VUgras{zaharniță} \textsuperscript{(43)}, va curăța
\VUgras{cleștișorul de zahăr} \textsuperscript{(44)} și va duce \VUgras{peria de
haine} \textsuperscript{(46)}.

%%K t06-c1-10-1 p
10. --- Partea dreaptă a camerei îi va da mai puțină treabă, fiindcă
\VUgras{canapeaua} \textsuperscript{(47)} și \VUgras{perna} ei
\textsuperscript{(48)} sunt la locul lor, iar fiindcă nu este frig, nu s-a atins
nimic nici din garnitura \VUgras{căminului} \textsuperscript{(49)}.
\VUgras{Lopățica} \textsuperscript{(50)}, \VUgras{cleștele}
\textsuperscript{(51)}, \VUgras{caprele} \textsuperscript{(55)} și
\VUgras{grătarul pentru cenușă} \textsuperscript{(56)} sunt curate;
\VUgras{paravanul căminului} \textsuperscript{(54)} nu este mișcat, iar
\VUgras{lada de lemne} \textsuperscript{(52)} este plină de \VUgras{lemne}
\textsuperscript{(53)}.

%%K t06-noto-1 noto
\VUnotes{143.2mm}{%
(*) Babo = copil mic; engl. \textit{baby}, fr. \textit{bébé}, ital.
\textit{bambino}.%
}

%%K t06-noto-2 noto
\VUnotes{143.2mm}{%
(*) Lambrequino = fr. \textit{lambrequin}, sp. \textit{lambrequines}, port.
\textit{lambrequins}, și chiar germ. și engl. \textit{lambrequins} --- adică și
nemțește, și englezește, și franțuzește, și portugheze, și spaniolește.%
}

%%K t06-c1-11-1 p
11. --- Va mai trebui să scuture \VUgras{ceasul de cămin} \textsuperscript{(57)},
\VUgras{sfeșnicele} \textsuperscript{(58)} și \VUgras{tăvițele pentru mărunțișuri}
\textsuperscript{(59)}, și apoi să șteargă \VUgras{oglinda}
\textsuperscript{(60)}.

%%K t06-c1-12-1 p
12. --- Cât despre \VUgras{leagănul} \textsuperscript{(61)} copilașului (babo (*)),
este gata.

%%K t06-tit-3 sub
\VUsaut{5.0mm}
\VUpk{10.2pt}{40}{Salonul.}
\VUsaut{3.0mm}

%%K t06-c2-01-1 p
1. --- Iată și salonul. Este o \VUgras{cameră} foarte frumoasă. Căminul din colțul
drept este împodobit cu o \VUgras{statuetă} elegantă \textsuperscript{(1)} și cu
două \VUgras{vaze} frumoase \textsuperscript{(3)}, iar drapat cu un
\VUgras{lambrechin} de catifea \textsuperscript{(2)} (*).

%%K t06-c2-02-1 p
2. --- Ca scânteile din vatră să nu aprindă covorul, înaintea căminului au pus un
\VUgras{paravan} de aramă \textsuperscript{(4)}.

%%K t06-c2-03-1 p
3. --- Alături stă deschis un \VUgras{birou} elegant de damă \textsuperscript{(5)}.
Aici \VUgras{stăpâna casei} \textsuperscript{(23)} își scrie scrisorile.

%%K t06-c2-04-1 p
4. --- Fereastra este împodobită cu \VUgras{perdele de tul} \textsuperscript{(6)},
strânse cu \VUgras{șnururi} \textsuperscript{(8)} pe \VUgras{cârlige}
\textsuperscript{(7)}, și cu draperii grele de stofă cu \VUgras{lambrechin}
\textsuperscript{(9)} deasupra.

%%K t06-c2-05-1 p
5. --- În firida ferestrei o \VUgras{jardinieră} \textsuperscript{(11)} ține o
\VUgras{plantă} verde \textsuperscript{(10)}, un palmier minunat, ale cărui frunze
ajung aproape până la \VUgras{lampa cu petrol} \textsuperscript{(12)}.

%%K t06-c2-06-1 p
6. --- \VUgras{Abajurul} \textsuperscript{(13)} l-a făcut Maria,
\VUgras{domnișoara} fermecătoare \textsuperscript{(14)} de la \VUgras{pian}
\textsuperscript{(15)}. Șezând foarte dreaptă pe \VUgras{taburet}
\textsuperscript{(16)}, învață o bucată. Este foarte iscusită și foarte ordonată,
ceea ce se vede din \VUgras{dulăpiorul ei de note} \textsuperscript{(17)}: acolo
este rânduială deplină.

%%K t06-tit-4 sub
\VUsaut{5.0mm}
\VUpk{10.2pt}{40}{Neastâmpăratul.}
\VUsaut{3.0mm}

%%K t06-c2-07-1 p
7. --- Fratele ei Pavel caută o \VUgras{carte} \textsuperscript{(19)} în
\VUgras{biblioteca} \textsuperscript{(18)}. Este un \VUgras{băiat}
\textsuperscript{(20)} neastâmpărat și cu totul de nesuferit. Atârnat de dulap ca
să ajungă la o carte prea sus așezată, se primejduiește să răstoarne
\VUgras{bustul} \textsuperscript{(21)} care stă deasupra.

%%K t06-c2-08-1 p
8. --- Se folosește pentru această prostie de faptul că mama este prinsă în vorbă
cu o prietenă, o \VUgras{doamnă} \textsuperscript{(22)} venită să stea la ea
câteva zile. Mama șade pe \VUgras{divan} \textsuperscript{(24)} lângă
\VUgras{fotoliu} \textsuperscript{(25)}, iar prietena ei pe un \VUgras{scaun}
\textsuperscript{(27)}, din care ne este văzută numai \VUgras{spătarul}
\textsuperscript{(26)}.

%%K t06-c2-09-1 p
9. --- În \VUgras{rama} mare \textsuperscript{(30)} de pe perete este
\VUgras{portretul} \textsuperscript{(29)} unei rude. În \VUgras{albumul}
\textsuperscript{(32)} care stă pe masă sunt strânse \VUgras{fotografiile}
\textsuperscript{(33)} celorlalte rude. Copiii tocmai l-au răsfoit, fiindcă
\VUgras{scaunele} \textsuperscript{(31)} sunt încă adunate în jurul mesei.

%%K t06-c2-10-1 p
10. --- \VUgras{Vaza chinezească} de porțelan \textsuperscript{(34)}, de lângă
\VUgras{cuțitul de hârtie} \textsuperscript{(35)}, i-a fost dăruită Mariei de ziua
onomastică, iar \VUgras{paravanul} \textsuperscript{(36)} care ocupă colțul camerei
l-a adus din China unchiul ei.

%%K t06-c2-11-1 p
11. --- Însuflețit de \VUgras{tablourile} frumoase \textsuperscript{(37)} de pe
\VUgras{perete} \textsuperscript{(42)}, luminat de \VUgras{candelabrul} din tavan
\textsuperscript{(38)}, ale cărui \VUgras{becuri electrice}
\textsuperscript{(39)} luminează atât de vesel seara, acest salon pare foarte
plăcut. \VUgras{Covorul} moale \textsuperscript{(40)}, întins pe podea, și
\VUgras{soneria electrică} atât de la îndemână \textsuperscript{(41)} îl fac cu
totul primitor.

%%K t06-tit-5 sub
\VUsaut{3.6mm}
\VUpk{10.2pt}{40}{Sufrageria.}
\VUsaut{3.0mm}

%%K t06-c3-01-1 p
1. --- S-a sunat la masă. Toată familia, afară de Maria, s-a strâns în jurul mesei.
Sufrageria este plăcut încălzită de o \VUgras{sobă} de teracotă aleasă
\textsuperscript{(1)}, cu \VUgras{burlanul} ei subțire \textsuperscript{(2)}.

%%K t06-c3-02-1 p
2. --- În această cameră este puțină mobilă. De-a lungul peretelui, deasupra
\VUgras{taburetului} \textsuperscript{(4)}, atârnă o \VUgras{fântâniță de spălat}
împodobită \textsuperscript{(3)}. Alături stă \VUgras{bufetul}
\textsuperscript{(5)}, pe care se pun bucatele reci, desertul și vasele felurite
care nu sunt încă de trebuință.

%%K t06-c3-03-1 p
3. --- Astfel, pe raftul de sus vedem o \VUgras{găină friptă}
\textsuperscript{(7)}, \VUgras{peștele} \textsuperscript{(8)} și \VUgras{vaza}
\textsuperscript{(10)} cu \VUgras{fructe} \textsuperscript{(9)}. Mai jos stau
\VUgras{brânza} \textsuperscript{(11)}, \VUgras{pâinea} \textsuperscript{(12)} și
\VUgras{serviciul} \textsuperscript{(13)}, în care este tot ce trebuie ca să dregi
salata: untdelemn, oțet, \VUgras{sare} \textsuperscript{(14)} și \VUgras{piper}
\textsuperscript{(15)}. În sfârșit, pe raftul de jos stă o \VUgras{plăcintă}
\textsuperscript{(17)} lângă \VUgras{muștariera} \textsuperscript{(16)}.

%%K t06-c3-04-1 p
4. --- Ceasul din sufragerie, care se numește \VUgras{de perete}
\textsuperscript{(20)}, este de o formă foarte neobișnuită. Este prins într-o ramă
de lemn, în care se află și \VUgras{barometrul} \textsuperscript{(19)} și
\VUgras{termometrul} \textsuperscript{(18)}.

%%K t06-tit-6 sub
\VUsaut{4.6mm}
\VUpk{10.2pt}{40}{Cum se dă masa.}
\VUsaut{3.0mm}

%%K t06-c3-05-1 p
5. --- De jur împrejurul camerei pereții sunt îmbrăcați în \VUgras{lambriuri de
lemn} \textsuperscript{(21)} din stejar ceruit. O \VUgras{draperie} de stofă
\textsuperscript{(22)} închide destul de bine ușa care duce spre verandă. Acolo va
merge familia după masă să bea cafea. \VUgras{Ceștile} \textsuperscript{(24)} și
\VUgras{farfurioarele} \textsuperscript{(25)} sunt deja așezate pe
\VUgras{etajeră} \textsuperscript{(23)}.

%%K t06-c3-06-1 p
6. --- Deasupra mesei este prinsă de tavan o \VUgras{lampă atârnată}
\textsuperscript{(26)}; lumina ei foarte vie umple seara toată sufrageria.

%%K t06-c3-07-1 p
7. --- Să privim acum \VUgras{masa} însăși \textsuperscript{(27)}. Când familia a
intrat, masa era pusă. Pe \VUgras{fața de masă} \textsuperscript{(28)} au așezat
la fiecare loc o \VUgras{farfurie} \textsuperscript{(33)}, un \VUgras{șervet}
\textsuperscript{(29)}, o \VUgras{lingură} \textsuperscript{(34)}, o
\VUgras{furculiță} \textsuperscript{(35)}, un \VUgras{cuțit}
\textsuperscript{(36)} și un \VUgras{pahar} \textsuperscript{(32)}. Erau acolo și
\VUgras{sticle} \textsuperscript{(30)} pentru vin și o \VUgras{carafă}
\textsuperscript{(31)} pentru apă.

%%K t06-c3-08-1 p
8. --- \VUgras{Servitorul} \textsuperscript{(39)} a adus mai întâi \VUgras{supiera}
\textsuperscript{(37)}, în care supa încă aburește. Acum dă \VUgras{felul} întâi
\textsuperscript{(38)}, cel de carne. Apoi îi va servi pe toți cu cele pe care
le-am numit; și în sfârșit, după \VUgras{desert}, se va folosi de faptul că
stăpânii au trecut pe verandă ca să strângă masa și să facă din nou rânduială în
sufragerie.

%%K t06-c3-08-2 p
\textit{Observație.} --- \textit{Cuvintele obișnuite despre mâncare sunt date aici
numai în linii mari. Lista va fi întregită în cele două tabele «Hotelul» (nr. 13)
și «Piața» (nr. 15).}

%%K t06-tit-7 sub
\VUsaut{4.6mm}
\VUpk{10.2pt}{40}{Bucătăria.}
\VUsaut{3.0mm}

%%K t06-c4-01-1 p
1. --- În bucătărie, unde privim prin ușa întredeschisă, ne întâmpină o
neorânduială pitorească. Înaintea \VUgras{vetrei} \textsuperscript{(1)}, unde
trosnește \VUgras{focul} \textsuperscript{(3)}, este așezată \VUgras{mașina de
frigare} \textsuperscript{(5)}. O \VUgras{friptură} minunată \textsuperscript{(7)}
este trasă în \VUgras{frigare} \textsuperscript{(6)} și se coace.

%%K t06-c4-02-1 p
2. --- \VUgras{Foalele} \textsuperscript{(8)} și \VUgras{lopățica de cărbuni}
\textsuperscript{(9)}, de care s-au folosit adineauri, au rămas pe podea împreună
cu \VUgras{butucii} \textsuperscript{(10)}.

%%K t06-c4-03-1 p
3. --- Pe polița căminului este rânduială: \VUgras{felinarul}
\textsuperscript{(11)}, \VUgras{cutia de chibrituri} \textsuperscript{(13)} cu
\VUgras{chibriturile} \textsuperscript{(12)}, \VUgras{sfeșnicul}
\textsuperscript{(14)} și \VUgras{sfeșnicul de noapte} \textsuperscript{(15)} cu
\VUgras{lumânarea} sa \textsuperscript{(16)} --- toate la locul lor. În schimb
\VUgras{găleata mare de cărbuni} \textsuperscript{(18)}, plină de
\VUgras{cărbuni} \textsuperscript{(17)}, după care \VUgras{bucătăreasa}
\textsuperscript{(23)} tocmai a coborât în pivniță, a fost lăsată în mijlocul
bucătăriei.

%%K t06-c4-04-1 p
4. --- Adevărul este că biata femeie trebuie să se grăbească, pentru ca prânzul să
fie gata la vreme. Acum este la \VUgras{mașina de gătit} \textsuperscript{(19)} și
amestecă \VUgras{sosul} \textsuperscript{(24)}, căruia nu i se poate îngădui să se
ardă.

%%K t06-c4-05-1 p
5. --- În \VUgras{cuptor} \textsuperscript{(20)} se coace plăcinta pe care a
făcut-o copiilor pentru gustare. Deasupra mașinii atârnă \VUgras{polonicul}
\textsuperscript{(21)} și \VUgras{spumiera} \textsuperscript{(22)}.

%%K t06-tit-8 sub
\VUsaut{4.0mm}
\VUpk{10.2pt}{40}{Vasele de bucătărie.}
\VUsaut{3.0mm}

%%K t06-c4-06-1 p
6. --- \VUgras{Garnitura de cratițe} \textsuperscript{(25)}, care atârnă în fundul
camerei, este foarte curată. \VUgras{Cratițele} frumoase de aramă
\textsuperscript{(26)}, \VUgras{oala de lapte} \textsuperscript{(27)},
\VUgras{strecurătoarea} \textsuperscript{(28)} și \VUgras{răzătoarea}
\textsuperscript{(29)} sunt curățate cu grijă.

%%K t06-c4-07-1 p
7. --- \VUgras{Solnița} \textsuperscript{(30)} stă pe un polițar mic lângă
\VUgras{ceainic} \textsuperscript{(31)} și lângă \VUgras{damigeana de untdelemn}
\textsuperscript{(32)}. În \VUgras{sertar} \textsuperscript{(33)}, se vede treaba,
stau tacâmurile și cuțitele de bucătărie.

%%K t06-c4-08-1 p
8. --- \VUgras{Mătura} \textsuperscript{(34)} și \VUgras{pămătuful de pene}
\textsuperscript{(35)} stau în colț. Bucătăreasa tocmai s-a folosit de
\VUgras{butucul de tăiat} \textsuperscript{(36)}; l-a lăsat lângă fereastră.

%%K t06-c4-09-1 p
9. --- \VUgras{Ștergarul} \textsuperscript{(37)} atârnă pe perete lângă
\VUgras{pâlnie} \textsuperscript{(38)}. În această parte a bucătăriei se țin
\VUgras{grătarul} \textsuperscript{(39)}, \VUgras{satârul}
\textsuperscript{(40)} și \VUgras{fundul de tăiat} \textsuperscript{(41)}. Apoi,
deasupra satârului, pe \VUgras{poliță} \textsuperscript{(42)}, stau
\VUgras{oalele} \textsuperscript{(43)}, \VUgras{tuciul} \textsuperscript{(44)} cu
\VUgras{capacul} său \textsuperscript{(45)} și \VUgras{cazanul}
\textsuperscript{(46)}. \VUgras{Tigaia} \textsuperscript{(47)} atârnă alături.

%%K t06-c4-10-1 p
10. --- Numai \VUgras{robinetul} de aramă \textsuperscript{(48)} aruncă o sclipire
în acest colț. \VUgras{Chiuveta} \textsuperscript{(49)}, pe care stă un
\VUgras{ulcior} mare \textsuperscript{(50)}, trebuie să fie foarte la îndemână cu
dulăpiorul ei, unde se țin \VUgras{butoiașul de oțet} \textsuperscript{(51)} cu
\VUgras{cepușorul} său \textsuperscript{(52)}, \VUgras{lada de gunoi}
\textsuperscript{(53)} și \VUgras{vasele murdare} \textsuperscript{(54)}.

%%K t06-tit-9 sub
\VUsaut{4.0mm}
\VUpk{10.2pt}{40}{Ce se mai ține în bucătărie.}
\VUsaut{3.0mm}

%%K t06-c4-11-1 p
11. --- \VUgras{Albia} mare \textsuperscript{(55)}, în care se spală
\VUgras{zarzavatul} \textsuperscript{(58)}, și \VUgras{cazanul} de fontă
\textsuperscript{(56)}, care stă pe \VUgras{pardoseala de teracotă}
\textsuperscript{(57)}, țin de bună seamă tot de acest dulăpior.

%%K t06-c4-12-1 p
12. --- Mașini de gătit bucătăresei nu îi lipsesc, căci iată aici, pe colțul mesei,
stă un \VUgras{bec de gaz} \textsuperscript{(59)}, iar pe el se încălzește apă
într-o cratiță. \VUgras{Aburul} \textsuperscript{(60)} care iese din ea ne spune
că apa trebuie să fiarbă. Apa aceasta este de bună seamă pentru cafea, fiindcă
iată la celălalt capăt al mesei \VUgras{ibricul de cafea} \textsuperscript{(64)}
și \VUgras{râșnița de cafea} \textsuperscript{(63)}.

%%K t06-c4-13-1 p
13. --- Becul este legat de \VUgras{țeava de gaz} \textsuperscript{(62)} cu un
\VUgras{furtun de cauciuc} lung \textsuperscript{(61)}.

%%K t06-c4-14-1 p
14. --- \VUgras{Femeia cu ziua} \textsuperscript{(66)}, care vine să o ajute pe
bucătăreasă, pregătește zarzavatul pentru masă. Când va fi curățat și spălat, îl va
pune în \VUgras{cratiță} \textsuperscript{(65)} până când va veni vremea să fie
fiert.

%%K t06-tit-10 sub
\VUsaut{4.0mm}
{\centering\fontsize{10.2pt}{10.2pt}\selectfont\textls[40]{\scshape Baia.} \textit{(Camera de scăldat.)}\par}
\VUsaut{3.0mm}

%%K t06-c5-01-1 p
1. --- Mai rămâne să privim într-un singur loc, ca să cunoaștem camerele de
căpetenie ale casei: să vedem întocmirea băii. Nu va fi lung, fiindcă este mică, și
vom vedea repede că \VUgras{cada} \textsuperscript{(1)} are un \VUgras{cazan} bun
\textsuperscript{(2)}, că \VUgras{săpuniera} \textsuperscript{(3)} este la
îndemâna celui care se scaldă, iar \VUgras{cuierul de prosoape}
\textsuperscript{(5)} trebuie să ușureze uscarea \VUgras{prosoapelor}
\textsuperscript{(4)}.

%%K t06-c5-02-1 p
2. --- Pereții acestei camere sunt acoperiți cu \VUgras{plăci smălțuite}
\textsuperscript{(6)}, ceea ce a îngăduit să se pună \VUgras{dușul}
\textsuperscript{(7)}, fără teama că apa care stropește la folosire îi va strica. Un
\VUgras{lighean} mic de zinc \textsuperscript{(8)}, pentru băi de picioare,
întregește mobilierul.
\VUsaut{6.0mm}
\VUfilet{22mm}
